\documentclass[11pt,a4paper]{article}
\usepackage[utf8]{inputenc}
\usepackage[T1]{fontenc}
\usepackage{lmodern}
\usepackage[french]{babel}
\usepackage{amsmath,amssymb,mathtools}
\usepackage{icomma}
\usepackage[version=4]{mhchem}
\usepackage{siunitx}
\sisetup{output-decimal-marker={,},per-mode=symbol}
\usepackage{xcolor}
\usepackage{colortbl}
\usepackage{tikz}
\usepackage[most]{tcolorbox}
\usepackage{enumitem}
\usepackage{array}
\usepackage{booktabs}
\usepackage{fancyhdr}
\usepackage[hidelinks]{hyperref}
\usepackage{titlesec}
\titlespacing*{\section}{0pt}{6pt}{3pt}
\usetikzlibrary{arrows.meta,positioning,calc,shapes.geometric,fit,backgrounds}
\usepackage[a4paper,margin=2.0cm,headheight=26pt,headsep=8pt]{geometry}

\definecolor{primary}{RGB}{146,95,161}
\definecolor{primarydark}{RGB}{92,55,108}
\definecolor{defcol}{RGB}{0,114,178}
\definecolor{thmcol}{RGB}{0,140,120}
\definecolor{formulacol}{RGB}{198,93,9}
\definecolor{remarkcol}{RGB}{120,94,200}
\definecolor{attncol}{RGB}{200,40,55}
\definecolor{excol}{RGB}{0,135,90}
\definecolor{methcol}{RGB}{90,90,100}
\definecolor{acidecol}{RGB}{205,60,45}
\definecolor{basecol}{RGB}{40,110,200}
\definecolor{protcol}{RGB}{240,150,30}
\definecolor{eaucol}{RGB}{0,160,190}
\definecolor{neutcol}{RGB}{90,170,90}

\tcbset{boxbase/.style={enhanced,breakable,boxrule=0.9pt,arc=2.5pt,left=3mm,right=3mm,top=2.3mm,bottom=2.3mm,fonttitle=\bfseries,coltitle=white,toptitle=0.6mm,bottomtitle=0.6mm}}
\newtcolorbox{definition}[1][]{boxbase,colback=defcol!5,colframe=defcol,title={\textsc{Définition}\ifx\relax#1\relax\relax\else\textmd{ : \itshape #1}\fi}}
\newtcolorbox{regle}[1][]{boxbase,colback=thmcol!6,colframe=thmcol,title={\textsc{Règle}\ifx\relax#1\relax\relax\else\textmd{ : \itshape #1}\fi}}
\newtcolorbox{formule}[1][]{boxbase,colback=formulacol!6,colframe=formulacol,title={\textsc{Formule}\ifx\relax#1\relax\relax\else\textmd{ : \itshape #1}\fi}}
\newtcolorbox{remarque}[1][]{boxbase,colback=remarkcol!6,colframe=remarkcol,title={\textsc{Remarque}\ifx\relax#1\relax\relax\else\textmd{ : \itshape #1}\fi}}
\newtcolorbox{attention}[1][]{boxbase,colback=attncol!6,colframe=attncol,title={\textsc{Attention}\ifx\relax#1\relax\relax\else\textmd{ : \itshape #1}\fi}}
\newtcolorbox{exemple}[1][]{boxbase,colback=excol!5,colframe=excol,title={\textsc{Exemple}\ifx\relax#1\relax\relax\else\textmd{ : \itshape #1}\fi}}
\newtcolorbox{methode}[1][]{boxbase,colback=methcol!7,colframe=methcol,sharp corners,title={\textsc{Méthode}\ifx\relax#1\relax\relax\else\textmd{ : \itshape #1}\fi}}

\pagestyle{fancy}
\fancyhf{}
\fancyhead[L]{\small\textcolor{primarydark}{\textbf{Fiche 10} — Thème 1 : couples acide/base}}
\fancyhead[R]{\small\textcolor{primarydark}{\textsc{Tuto}}}
\fancyfoot[C]{\small\thepage}
\renewcommand{\headrulewidth}{0.8pt}
\renewcommand{\headrule}{\hbox to\headwidth{\color{primary}\leaders\hrule height \headrulewidth\hfill}}

\setlength{\parindent}{0pt}
\setlength{\parskip}{3pt}

\begin{document}

\begin{center}
{\LARGE\bfseries\color{primarydark} Thème 1 --- Définitions, couples acide/base conjugués et ampholytes}\\[3pt]
{\color{primary}\itshape Fiche tuto à lire avant les exercices}
\end{center}

\section*{1. La seule idée à retenir : un proton \ce{H+} qui change de mains}

Toute la chimie acide-base de Brønsted-Lowry tient dans \textbf{un transfert de proton \ce{H+}} (un simple
noyau d'hydrogène). On a besoin de \emph{deux} définitions symétriques :

\begin{definition}[acide et base selon Brønsted-Lowry]
Un \textbf{acide} est une espèce capable de \textbf{céder} (donner) un proton \ce{H+}.\\
Une \textbf{base} est une espèce capable de \textbf{capter} (accepter) un proton \ce{H+}.
\end{definition}

Un acide ne peut donner son \ce{H+} que s'il y a une base pour le recevoir : acide et base fonctionnent
\textbf{toujours par deux}. Attention, le proton \ce{H+} n'existe jamais seul dans l'eau : il est solvaté sous
la forme \ce{H3O+} (ion oxonium). On écrit donc indifféremment \ce{H+} ou \ce{H3O+}.

\section*{2. Le couple acide/base conjugués : la mécanique centrale du thème}

Quand un acide \ce{HA} cède son proton, il ne disparaît pas : il se transforme en une base \ce{A-}, sa
\textbf{base conjuguée}. L'ensemble forme un \textbf{couple} noté \ce{HA}/\ce{A-}.

\begin{regle}[construire le partenaire d'un couple --- le geste à automatiser]
\begin{itemize}[leftmargin=1.5em,topsep=1pt,itemsep=1pt]
  \item \textbf{Base conjuguée d'un acide} : on \textbf{enlève un \ce{H+}} $\Rightarrow$ un \ce{H} en moins
        \emph{et} une charge en moins (\,$+1 \to 0$, ou $0 \to -1$, ou $-1 \to -2$\,).
  \item \textbf{Acide conjugué d'une base} : on \textbf{ajoute un \ce{H+}} $\Rightarrow$ un \ce{H} en plus
        \emph{et} une charge en plus.
\end{itemize}
La demi-équation d'un couple s'écrit toujours : $\ce{HA <=> A- + H+}$, ou en milieu aqueux
$\ce{HA + H2O <=> A- + H3O+}$.
\end{regle}

\begin{exemple}[on enlève / on ajoute un \ce{H+}, rien d'autre]
\begin{center}\small
\begin{tabular}{lcl@{\qquad}lcl}
\toprule
Acide & $\to$ & Base conjuguée & Base & $\to$ & Acide conjugué\\
\midrule
\ce{HCl}        & $-\ce{H+}$ & \ce{Cl-}            & \ce{NH3}          & $+\ce{H+}$ & \ce{NH4+}\\
\ce{CH3COOH}    & $-\ce{H+}$ & \ce{CH3COO-}        & \ce{OH-}          & $+\ce{H+}$ & \ce{H2O}\\
\ce{NH4+}       & $-\ce{H+}$ & \ce{NH3}            & \ce{CO3^{2-}}     & $+\ce{H+}$ & \ce{HCO3-}\\
\ce{H2CO3}      & $-\ce{H+}$ & \ce{HCO3-}          & \ce{H2O}          & $+\ce{H+}$ & \ce{H3O+}\\
\bottomrule
\end{tabular}
\end{center}
Remarquez comme la charge suit toujours le proton : $\ce{NH4+}\ (+1) \to \ce{NH3}\ (0)$,
$\ce{CO3^{2-}}\ (-2) \to \ce{HCO3-}\ (-1)$.
\end{exemple}

\begin{attention}[deux bornes que l'eau ne dépasse pas]
En solution aqueuse, \ce{H3O+} est l'acide le plus fort qui existe (son « acide conjugué » \ce{H4O^{2+}}
n'existe pas), et \ce{OH-} est la base la plus forte usuelle. Donc \textbf{\ce{H3O+} n'a pas d'acide conjugué
observable}, et la base conjuguée de \ce{OH-} (l'ion \ce{O^{2-}}) n'existe pas dans l'eau.
\end{attention}

\section*{3. Les ampholytes (espèces amphiprotiques)}

\begin{definition}[ampholyte]
Un \textbf{ampholyte} est une espèce qui peut jouer \textbf{soit le rôle d'acide, soit le rôle de base},
selon son partenaire. Il en existe dès qu'une espèce possède à la fois un \ce{H} cessible \emph{et} un site
capable de capter un \ce{H+}. Un ampholyte appartient donc à \textbf{deux couples différents}.
\end{definition}

\begin{exemple}[\ce{H2O} et \ce{HCO3-}, les ampholytes stars]
\textbf{L'eau} est l'ampholyte le plus important : acide du couple \ce{H2O}/\ce{OH-}, base du couple
\ce{H3O+}/\ce{H2O}.\\[2pt]
\textbf{L'hydrogénocarbonate \ce{HCO3-}} :
\[ \underbrace{\ce{HCO3- + H2O <=> CO3^{2-} + H3O+}}_{\text{ici \ce{HCO3-} est l'\textbf{acide} du couple }\ce{HCO3-}/\ce{CO3^{2-}}} \qquad
   \underbrace{\ce{HCO3- + H3O+ <=> H2CO3 + H2O}}_{\text{ici \ce{HCO3-} est la \textbf{base} du couple }\ce{H2CO3}/\ce{HCO3-}} \]
On repère un ampholyte à sa position \emph{au milieu} d'une chaîne de déprotonations, par exemple pour le
triacide \ce{H3PO4} : \ce{H3PO4} $\to$ \ce{H2PO4-} $\to$ \ce{HPO4^{2-}} $\to$ \ce{PO4^{3-}}. Les espèces
intermédiaires \ce{H2PO4-} et \ce{HPO4^{2-}} sont des \textbf{ampholytes}.
\end{exemple}

\section*{4. Les deux autres théories (culture utile)}

\textbf{Arrhenius} (limitée à l'eau) : un acide libère des \ce{H3O+}, une base libère des \ce{OH-}.
\quad\textbf{Lewis} (la plus générale) : un \textbf{acide de Lewis} \emph{accepte} un doublet électronique
(atome à orbitale vide : \ce{H+}, \ce{Ag+}, \ce{BF3}) ; une \textbf{base de Lewis} \emph{donne} un doublet
(\ce{NH3}, \ce{H2O}). Tout acide de Brønsted est un acide de Lewis, mais l'inverse est faux.

\begin{methode}[le réflexe du thème]
Devant une espèce : peut-elle \emph{donner} un \ce{H+} $\Rightarrow$ acide ; \emph{capter} un \ce{H+}
$\Rightarrow$ base ; \emph{les deux} $\Rightarrow$ ampholyte. Pour trouver son partenaire, ne touchez qu'\textbf{à
un seul \ce{H} et une seule charge}.
\end{methode}

\end{document}
